\section{Adjoint functors}
\label{section-adjoint}

\begin{definition}
\label{definition-adjoint}
Let $\mathcal{C}$, $\mathcal{D}$ be categories.
Let $u : \mathcal{C} \to \mathcal{D}$ and
$v : \mathcal{D} \to \mathcal{C}$ be functors.
We say that $u$ is a {\it left adjoint} of $v$, or that
$v$ is a {\it right adjoint} to $u$ if there are bijections
$$
\Mor_\mathcal{D}(u(X), Y)
\longrightarrow
\Mor_\mathcal{C}(X, v(Y))
$$
functorial in $X \in \Ob(\mathcal{C})$, and
$Y \in \Ob(\mathcal{D})$.
\end{definition}

\noindent
In other words, this means that there is a {\it given} isomorphism of functors
$\mathcal{C}^{opp} \times \mathcal{D} \to \textit{Sets}$ from
$\Mor_\mathcal{D}(u(-), -)$ to $\Mor_\mathcal{C}(-, v(-))$. For any object
$X$ of $\mathcal{C}$ we obtain a morphism $X \to v(u(X))$ corresponding to
$\text{id}_{u(X)}$. Similarly, for any object $Y$ of $\mathcal{D}$ we obtain
a morphism $u(v(Y)) \to Y$ corresponding to $\text{id}_{v(Y)}$.
These maps are called the {\it adjunction maps}. The adjunction maps
are functorial in $X$ and $Y$, hence we obtain morphisms of functors
$$
\eta : \text{id}_\mathcal{C} \to v \circ u\quad (\text{unit})
\quad\text{and}\quad
\epsilon : u \circ v \to \text{id}_\mathcal{D}\quad (\text{counit}).
$$
Moreover, if $\alpha : u(X) \to Y$
and $\beta : X \to v(Y)$ are morphisms, then the following are equivalent
\begin{enumerate}
\item $\alpha$ and $\beta$ correspond to each other via the
bijection of the definition,
\item $\beta$ is the composition $X \to v(u(X)) \xrightarrow{v(\alpha)} v(Y)$,
and
\item $\alpha$ is the composition $u(X) \xrightarrow{u(\beta)} u(v(Y)) \to Y$.
\end{enumerate}
In this way one can reformulate the notion of adjoint functors in terms
of adjunction maps.

\begin{lemma}
\label{lemma-adjoint-exists}
Let $u : \mathcal{C} \to \mathcal{D}$ be a functor between categories.
If for each $y \in \Ob(\mathcal{D})$ the functor
$x \mapsto \Mor_\mathcal{D}(u(x), y)$ is representable, then
$u$ has a right adjoint.
\end{lemma}

\begin{proof}
For each $y$ choose an object $v(y)$ and an isomorphism
$\Mor_\mathcal{C}(-, v(y)) \to \Mor_\mathcal{D}(u(-), y)$
of functors. By Yoneda's lemma (Lemma \ref{lemma-yoneda})
for any morphism $g : y \to y'$ the transformation of functors
$$
\Mor_\mathcal{C}(-, v(y)) \to \Mor_\mathcal{D}(u(-), y) \to
\Mor_\mathcal{D}(u(-), y') \to \Mor_\mathcal{C}(-, v(y'))
$$
corresponds to a unique morphism $v(g) : v(y) \to v(y')$.
We omit the verification that $v$ is a functor and that
it is right adjoint to $u$.
\end{proof}

\begin{lemma}
\label{lemma-left-adjoint-composed-fully-faithful}
\begin{reference}
Bhargav Bhatt, private communication.
\end{reference}
Let $u$ be a left adjoint to $v$ as in Definition \ref{definition-adjoint}.
\begin{enumerate}
\item If $v \circ u$ is fully faithful, then $u$ is fully faithful.
\item If $u \circ v$ is fully faithful, then $v$ is fully faithful.
\end{enumerate}
\end{lemma}

\begin{proof}
Proof of (2). Assume $u \circ v$ is fully faithful.
Say we have $X$, $Y$ in $\mathcal{D}$.
Then the natural composite map
$$
\Mor(X,Y) \to \Mor(v(X),v(Y)) \to \Mor(u(v(X)), u(v(Y)))
$$
is a bijection, so $v$ is at least faithful. To show full faithfulness,
we must show that the second map above is injective.
But the adjunction between $u$ and $v$ says that
$$
\Mor(v(X), v(Y)) \to \Mor(u(v(X)), u(v(Y))) \to \Mor(u(v(X)), Y)
$$
is a bijection, where the first map is natural one and
the second map comes from the counit $u(v(Y)) \to Y$ of the adjunction.
So this says that
$\Mor(v(X), v(Y)) \to \Mor(u(v(X)), u(v(Y)))$
is also injective, as wanted. The proof of (1) is dual to this.
\end{proof}

\begin{lemma}
\label{lemma-adjoint-fully-faithful}
Let $u$ be a left adjoint to $v$ as in Definition \ref{definition-adjoint}.
Then
\begin{enumerate}
\item $u$ is fully faithful $\Leftrightarrow$ $\text{id} \cong v \circ u$
$\Leftrightarrow$ $\eta : \text{id} \to v \circ u$ is an isomorphism,
\item $v$ is fully faithful $\Leftrightarrow$
$u \circ v \cong \text{id}$ $\Leftrightarrow$
$\epsilon : u \circ v \to \text{id}$ is an isomorphism.
\end{enumerate}
\end{lemma}

\begin{proof}
Proof of (1).
Assume $u$ is fully faithful. We will show $\eta_X : X \to v(u(X))$
is an isomorphism. Let $X' \to v(u(X))$ be any morphism.
By adjointness this corresponds to a morphism $u(X') \to u(X)$. By fully
faithfulness of $u$ this corresponds to a unique morphism $X' \to X$.
Thus we see that post-composing by $\eta_X$ defines a bijection
$\Mor(X', X) \to \Mor(X', v(u(X)))$. Hence $\eta_X$ is an isomorphism.
If there exists an isomorphism $\text{id} \cong v \circ u$ of functors,
then $v \circ u$ is fully faithful. By
Lemma \ref{lemma-left-adjoint-composed-fully-faithful} we see
that $u$ is fully faithful. By the above this implies $\eta$
is an isomorphism. Thus all $3$ conditions are equivalent (and these
conditions are also equivalent to $v \circ u$ being fully faithful).

\medskip\noindent
Part (2) is dual to part (1).
\end{proof}

\begin{lemma}
\label{lemma-adjoint-exact}
Let $u$ be a left adjoint to $v$ as in Definition \ref{definition-adjoint}.
\begin{enumerate}
\item Suppose that $M : \mathcal{I} \to \mathcal{C}$ is a diagram,
and suppose that $\colim_\mathcal{I} M$ exists in
$\mathcal{C}$. Then $u(\colim_\mathcal{I} M) =
\colim_\mathcal{I} u \circ M$. In other words,
$u$ commutes with (representable) colimits.
\item Suppose that $M : \mathcal{I} \to \mathcal{D}$ is a diagram,
and suppose that $\lim_\mathcal{I} M$ exists in
$\mathcal{D}$. Then $v(\lim_\mathcal{I} M) =
\lim_\mathcal{I} v \circ M$. In other words $v$ commutes
with representable limits.
\end{enumerate}
\end{lemma}

\begin{proof}
A morphism from a colimit into an object is the same as a compatible
system of morphisms from the constituents of the limit into the
object, see Remark \ref{remark-limit-colim}. So
$$
\begin{matrix}
\Mor_\mathcal{D}(u(\colim_{i \in \mathcal{I}} M_i), Y) &
= & \Mor_\mathcal{C}(\colim_{i \in \mathcal{I}} M_i, v(Y)) \\
& = &
\lim_{i \in \mathcal{I}^{opp}}
\Mor_\mathcal{C}(M_i, v(Y)) \\
& = &
\lim_{i \in \mathcal{I}^{opp}}
\Mor_\mathcal{D}(u(M_i), Y)
\end{matrix}
$$
proves that $u(\colim_{i \in \mathcal{I}} M_i)$ is
the colimit we are looking for.
A similar argument works for the other statement.
\end{proof}

\begin{lemma}
\label{lemma-exact-adjoint}
Let $u$ be a left adjoint of $v$ as in Definition \ref{definition-adjoint}.
\begin{enumerate}
\item If $\mathcal{C}$ has finite colimits, then $u$ is right exact.
\item If $\mathcal{D}$ has finite limits, then $v$ is left exact.
\end{enumerate}
\end{lemma}

\begin{proof}
Obvious from the definitions and Lemma \ref{lemma-adjoint-exact}.
\end{proof}

\begin{lemma}
\label{lemma-unit-counit-relations}
Let $u : \mathcal{C} \to \mathcal{D}$ be a left adjoint to the functor
$v : \mathcal{D} \to \mathcal{C}$. Let $\eta_X : X \to v(u(X))$ be the unit
and $\epsilon_Y : u(v(Y)) \to Y$ be the counit. Then
$$
u(X) \xrightarrow{u(\eta_X)} u(v(u(X))
\xrightarrow{\epsilon_{u(X)}} u(X)
\quad\text{and}\quad
v(Y) \xrightarrow{\eta_{v(Y)}} v(u(v(Y))) \xrightarrow{v(\epsilon_Y)}
v(Y)
$$
are the identity morphisms.
\end{lemma}

\begin{proof}
Omitted.
\end{proof}

\begin{lemma}
\label{lemma-transformation-between-functors-and-adjoints}
Let $u_1, u_2 : \mathcal{C} \to \mathcal{D}$ be functors with right
adjoints $v_1, v_2 : \mathcal{D} \to \mathcal{C}$. Let $\beta : u_2 \to u_1$
be a transformation of functors. Let $\beta^\vee : v_1 \to v_2$ be
the corresponding transformation of adjoint functors. Then
$$
\xymatrix{
u_2 \circ v_1 \ar[r]_\beta \ar[d]_{\beta^\vee} &
u_1 \circ v_1 \ar[d] \\
u_2 \circ v_2 \ar[r] & \text{id}
}
$$
is commutative where the unlabeled arrows are the counit transformations.
\end{lemma}

\begin{proof}
This is true because $\beta^\vee_D : v_1D \to v_2D$ is the unique
morphism such that the induced maps $\Mor(C, v_1D) \to \Mor(C, v_2D)$
is the map $\Mor(u_1C, D) \to \Mor(u_2C, D)$ induced by
$\beta_C : u_2C \to u_1C$. Namely, this means the map
$$
\Mor(u_1 v_1 D, D') \to \Mor(u_2 v_1 D, D')
$$
induced by $\beta_{v_1 D}$ is the same as the map
$$
\Mor(v_1 D, v_1 D') \to \Mor(v_1 D, v_2 D')
$$
induced by $\beta^\vee_{D'}$. Taking $D' = D$ we find that the counit
$u_1 v_1 D \to D$ precomposed by $\beta_{v_1D}$ corresponds to $\beta^\vee_D$
under adjunction. This exactly means that the diagram commutes when
evaluated on $D$.
\end{proof}

\begin{lemma}
\label{lemma-compose-counits}
Let $\mathcal{A}$, $\mathcal{B}$, and $\mathcal{C}$ be categories.
Let $v : \mathcal{A} \to \mathcal{B}$ and
$v' : \mathcal{B} \to \mathcal{C}$ be functors
with left adjoints $u$ and $u'$ respectively. Then
\begin{enumerate}
\item The functor $v'' = v' \circ v$ has a left adjoint equal to
$u'' = u \circ u'$.
\item Given $X$ in $\mathcal{A}$ we have
\begin{equation}
\label{equation-compose-counits}
\epsilon_X^v \circ u(\epsilon^{v'}_{v(X)}) = \epsilon^{v''}_X :
u''(v''(X)) \to X
\end{equation}
Where $\epsilon$ is the counit of the adjunctions.
\end{enumerate}
\end{lemma}

\begin{proof}
Let us unwind the formula in (2) because this will also immediately
prove (1). First, the counit of the adjunctions for the pairs
$(u, v)$ and $(u', v')$ are maps
$\epsilon_X^v : u(v(X)) \to X$ and
$\epsilon_Y^{v'} : u'(v'(Y)) \to Y$, see discussion following
Definition \ref{definition-adjoint}.
With $u''$ and $v''$ as in (1) we unwind everything
$$
u''(v''(X)) = u(u'(v'(v(X)))) \xrightarrow{u(\epsilon_{v(X)}^{v'})}
u(v(X)) \xrightarrow{\epsilon_X^v} X
$$
to get the map on the left hand side of (\ref{equation-compose-counits}).
Let us denote this by $\epsilon_X^{v''}$ for now.
To see that this is the counit of an adjoint pair
$(u'', v'')$ we have to show that given $Z$ in $\mathcal{C}$
the rule that sends a morphism $\beta : Z \to v''(X)$
to $\alpha = \epsilon_X^{v''} \circ u''(\beta) : u''(Z) \to X$
is a bijection on sets of morphisms.
This is true because, this is the composition of the
rule sending $\beta$ to $\epsilon_{v(X)}^{v'} \circ u'(\beta)$
which is a bijection by assumption on $(u', v')$ and then
sending this to
$\epsilon_X^v \circ u(\epsilon_{v(X)}^{v'} \circ u'(\beta))$
which is a bijection by assumption on $(u, v)$.
\end{proof}
